% =====================================================================
% Wave borcherds_chain_level — Agent BL-7
% MASTER THEOREM: chain-level identity
%
%   kappa_BKM(Phi_N) = c_N(0)/2
%                   = STr_{mod.op}[B^{E_3}(Phi^{FA}_3(C_N))]
%
% for C_N a CY_3 category indexed by N in {1,2,3,4,6} (CHL ladder;
% Scope A) or by Gritsenko–Clery 8-form index (Scope B).
%
% Upstream prior-wave inputs synthesised:
%   Agent  1 (KT E_3-formality, GRT_1(Q)-torsor on Stage-1)
%   Agent  2 (factorisation envelope of Phi^{FA}_3)
%   Agent  6 (Kapranov L_infty on T_X[-1], three Atiyah cocycles)
%   Agent  7 (Dunn–Lurie E_3 ≃ E_1 ⊗ E_2 chain-level)
%   Agent  9 (Fulton–MacPherson configuration operad)
%   Agent 11 (E_3 Koszul self-duality up to shift)
%   Agent 13 (obstruction theory; O_2 Mumford collapse)
%   Agent 14 (Cech–Ran descent; super-Yangian emergence)
%   Agent 15 (modular Koszul bridge; Torelli–Kuga–Satake)
%
% Standalone notes-lane file. Manuscript inscription goes into
%   chapters/examples/cy_d_kappa_stratification.tex
% at label
%   \label{thm:bl-master-kappa-bkm-chain-level}
%
% Primary references:
%   Getzler–Kapranov "Modular operads" Compositio Math. 110 (1998) 65–126.
%   Borcherds "Automorphic forms on O_{s+2,2}(R) and infinite products"
%     Invent. Math. 120 (1995) 161–213.
%   Borcherds "Automorphic forms with singularities on Grassmannians"
%     Invent. Math. 132 (1998) 491–562.
%   Gritsenko–Nikulin "Automorphic forms and Lorentzian Kac–Moody
%     algebras II" alg-geom/9504006 (1995).
%   Gritsenko "Arithmetical lifting and its applications" in
%     Algebraic Geometry Santa Cruz 1995 (1999).
%   Gritsenko–Clery "Siegel modular forms of genus 2 spanned by
%     theta series" arXiv:1305.4884 (2013).
%   Ayala–Francis "Poincaré/Koszul duality" arXiv:1409.2478 (2014).
%   Costello–Gwilliam "Factorization Algebras in QFT" vol. II (2021).
%   Willwacher "Kontsevich's graph complex and GRT_1" Invent. Math.
%     200 (2015) 671–760.
%   Bruinier "Borcherds products on O(2,l) and Chern classes of
%     Heegner divisors" Lecture Notes in Mathematics 1780 (2002).
%   Lorgat "Automorphic corrections for the GKM superalgebra
%     g_{Delta_5}" arXiv:2004.09030 (2020).
%   Costello–Francis–Gwilliam "Topological Chern–Simons" (2026).
%
% =====================================================================

\documentclass[11pt]{amsart}
\usepackage{amsmath,amssymb,amsthm,mathtools,hyperref,enumitem}
\usepackage[margin=1in]{geometry}
\newtheorem{theorem}{Theorem}[section]
\newtheorem{proposition}[theorem]{Proposition}
\newtheorem{lemma}[theorem]{Lemma}
\newtheorem{corollary}[theorem]{Corollary}
\newtheorem{conjecture}[theorem]{Conjecture}
\newtheorem{definition}[theorem]{Definition}
\newtheorem{remark}[theorem]{Remark}
\newtheorem{attack}[theorem]{Attack}
\newtheorem{surviving}[theorem]{Surviving core}
\newtheorem{heal}[theorem]{Heal}
\newtheorem{ghost}[theorem]{Ghost theorem}

\providecommand{\C}{\mathbb{C}}
\providecommand{\R}{\mathbb{R}}
\providecommand{\Z}{\mathbb{Z}}
\providecommand{\Q}{\mathbb{Q}}
\providecommand{\HH}{\mathrm{HH}}
\providecommand{\CH}{\mathrm{CH}}
\providecommand{\CC}{\mathrm{CC}}
\providecommand{\Ran}{\mathrm{Ran}}
\providecommand{\cC}{\mathcal{C}}
\providecommand{\cA}{\mathcal{A}}
\providecommand{\cB}{\mathcal{B}}
\providecommand{\cD}{\mathcal{D}}
\providecommand{\cE}{\mathcal{E}}
\providecommand{\cM}{\mathcal{M}}
\providecommand{\cO}{\mathcal{O}}
\providecommand{\cR}{\mathcal{R}}
\providecommand{\cS}{\mathcal{S}}
\providecommand{\cT}{\mathcal{T}}
\providecommand{\fg}{\mathfrak{g}}
\providecommand{\fh}{\mathfrak{h}}
\providecommand{\sp}{\mathrm{sp}}
\providecommand{\Sp}{\mathrm{Sp}}
\providecommand{\Mp}{\mathrm{Mp}}
\providecommand{\Mpt}{\widetilde{\mathrm{Mp}}}
\providecommand{\GRT}{\mathrm{GRT}}
\providecommand{\PhiFA}{\Phi^{\mathrm{FA}}}
\providecommand{\SpCh}{\mathrm{Sp}^{\mathrm{ch}}}
\providecommand{\Mbar}{\overline{\mathcal{M}}}
\providecommand{\Abar}{\overline{\mathcal{A}}}
\providecommand{\Perf}{\mathrm{Perf}}
\providecommand{\Coh}{\mathrm{Coh}}
\providecommand{\At}{\mathrm{At}}
\providecommand{\Ger}{\mathrm{Ger}}
\providecommand{\KS}{\mathrm{KS}}
\providecommand{\FM}{\mathrm{FM}}
\providecommand{\BV}{\mathrm{BV}}
\providecommand{\MO}{\mathrm{MO}}
\providecommand{\Hodge}{\mathrm{Hodge}}
\providecommand{\Borch}{\mathrm{Borch}}
\providecommand{\CHL}{\mathrm{CHL}}
\providecommand{\Ber}{\mathrm{Ber}}
\providecommand{\Hdg}{\mathrm{Hdg}}
\providecommand{\STr}{\mathrm{STr}}
\providecommand{\modop}{\mathrm{mod.op}}
\providecommand{\kBKM}{\kappa_{\mathrm{BKM}}}
\providecommand{\kch}{\kappa_{\mathrm{ch}}}
\providecommand{\kcat}{\kappa_{\mathrm{cat}}}
\providecommand{\kfib}{\kappa_{\mathrm{fiber}}}
\providecommand{\End}{\mathrm{End}}
\providecommand{\Ext}{\mathrm{Ext}}
\providecommand{\Aut}{\mathrm{Aut}}
\providecommand{\Hom}{\mathrm{Hom}}
\providecommand{\Tor}{\mathrm{Tor}}
\providecommand{\id}{\mathrm{id}}

\begin{document}

\title{The chain-level master identity\\
$\kBKM(\Phi_N) = c_N(0)/2 = \STr_{\modop}[B^{E_3}(\PhiFA_3(\cC_N))]$}
\author{Agent BL-7, wave borcherds\_chain\_level, Vol III}
\date{April 2026}
\maketitle

\begin{abstract}
For every $N$ in either the CHL ladder $\{1,2,3,4,6\}$ (Scope A) or
the Gritsenko--Cl\'ery $8$-form index $\{1,2,3,4,5,6,1/2,3/2\}$
(Scope B), the Borcherds weight $\kBKM(\Phi_N)$ of the paramodular
Borcherds product $\Phi_N$ equals a chain-level modular-operadic
supertrace of the Koszul dual of the Stage-$1$ $E_3$-factorisation
algebra attached to a CY$_3$ category $\cC_N$:
\[
 \kBKM(\Phi_N) \;=\; \frac{c_N(0)}{2}
 \;=\; \STr_{\modop}\bigl[B^{E_3}\bigl(\PhiFA_3(\cC_N)\bigr)\bigr].
\]
We construct the modular-operadic supertrace from the Getzler--Kapranov
$1998$ modular operad on $\overline{\mathcal M}_{g,n}$ via
Torelli--Kuga--Satake uniformisation of the K3 period domain, localise
to the Borcherds lift generating series at genus $g = 2$, and establish
three-path compatibility: the Hodge-supertrace path (Kapranov
$L_\infty$ on $T_X[-1]$), the BV-residue path (Costello--Li--Yagi
partition function at Heegner locus), and the MO-residue path
(Maulik--Okounkov $R$-matrix at UV gluing cocycle wall) all produce the
same chain-level numeric. The three quasi-isomorphisms assemble into a
canonical chain-level Siegel modular form on
$\mathrm{Sp}_4(\Z)\backslash \mathbb{H}_2$ (or on the metaplectic /
double-metaplectic cover for half- / quarter-integral weights of
Scope B), whose zeroth Fourier coefficient is $c_N(0)/2$.
GRT$_1(\Q)$-invariance of the numeric is established via the action
of $\GRT_1(\Q)$ on Stage-$1$ lifts and the modular operad. The
two-stage factorisation discipline is preserved: Stage-$2$
specialisation is compatible with the modular-operadic pairing because
the genus-tower curvature collapses via the Mumford boundary relation
$\partial^* \lambda_g^1 = \pi_1^* \lambda_{g_1}^1 + \pi_2^* \lambda_{g_2}^1$,
identifying the $\kch$-curving with the Hodge-boundary-descent datum
(Agent 13). The worked example $\cC_1 = D^b\Coh(K3 \times E)$ recovers
$\kBKM(\fg_{\Delta_5}) = 5$ via Lorgat $2020$'s explicit Borcherds
lift of $\phi_{0,1}$. The file runs five attack--heal cycles plus a
ghost-theorem harvest.
\end{abstract}

\tableofcontents

% =====================================================================
\section{Standing data and the master identity}
\label{sec:master-standing-data}
% =====================================================================

Fix once and for all.
\begin{itemize}
\item $k = \Q$ (rational coefficients; Drinfeld-associator discipline of
Agent $1$, \S~4).
\item $\cC_N \in \mathrm{CY}_3\text{-}\mathrm{Cat}^{\mathrm{dg}}$ a
smooth proper cyclic $A_\infty$-category of CY dimension $3$ over
$\Q$, with $\HH^0(\cC_N) = \Q$, indexed by either
  \begin{itemize}
  \item[(A)] $N \in \{1, 2, 3, 4, 6\}$, the \emph{CHL ladder}:
  $\cC_N = D^b \Coh^{g_N}(\mathrm{K3} \times E)$, the derived
  category of coherent sheaves equivariant under a
  symplectic automorphism $g_N \in M_{24}$ of order $N$; or
  \item[(B)] $N \in \{1, 2, 3, 4, 5, 6, 1/2, 3/2\}$, the
  \emph{Gritsenko--Cl\'ery $8$-form index}: $\cC_N$ the CY$_3$
  category associated to the paramodular level-$m$ host at integer or
  half-integer paramodular index $m = N$, with
  $\cC_1 = D^b\Coh(\mathrm{K3} \times E)$ at $N = 1$.
  \end{itemize}
\item $\PhiFA_3$ the Stage-$1$ two-stage-factorisation functor of
Proposition~\ref{prop:phifa-infty1-kernel} (working notes):
\[
 \PhiFA_3 \colon
 \mathrm{CY}_3\text{-}\mathrm{Cat}^{\mathrm{dg}}
 \;\longrightarrow\;
 E_3\text{-}\HolFA(X),
\]
canonical up to a contractible $\GRT_1(\Q)$-torsor of choices
(Agent $1$, Theorem~\ref{thm:form3-grt1-torsor}).
\item $B^{E_3}$ the $E_3$-Koszul-bar functor on $E_3$-algebras over
$\Q$, valued in $E_3$-coalgebras (Ayala--Francis $2014$
Theorem~$2.1$).
\item $\Phi_N$ the paramodular Borcherds product attached to $\cC_N$,
with zeroth Fourier coefficient $c_N(0)$ of the Borcherds-input
principal part.
\item $\kBKM(\Phi_N) = \mathrm{wt}(\Phi_N)$ the Borcherds weight,
equal to $c_N(0)/2$ by Borcherds $1998$ Theorem~$13.3$.
\item $\mathsf{MMG}$ the Getzler--Kapranov $1998$ modular operad with
$\mathsf{MMG}(g, n) = \overline{\mathcal M}_{g, n}$; $[\Mbar_{g, n}]
\in C_\bullet^{\FM\text{-}\log}(\Mbar_{g, n}; \Q)$ its chain-level
fundamental class (Agent $15$ Heal~I.$1$).
\end{itemize}

\begin{definition}[The master identity]
\label{def:master-identity}
The \emph{chain-level master identity} at scope (A) or (B) is the
equation
\begin{equation}
\label{eq:master-chain-level}
 \kBKM(\Phi_N) \;=\; \frac{c_N(0)}{2} \;=\;
 \STr_{\modop}\bigl[B^{E_3}\bigl(\PhiFA_3(\cC_N)\bigr)\bigr],
\end{equation}
interpreted as follows:
\begin{enumerate}[label=\textup{(\roman*)}]
\item the first equality is Borcherds' singular-theta weight formula
($1998$ Theorem~$13.3$), a statement about the Borcherds product
$\Phi_N$;
\item the second equality is the chain-level master theorem of this
file, identifying that weight as a modular-operadic supertrace of the
Koszul dual of the Stage-$1$ $E_3$-algebra attached to $\cC_N$;
\item $\STr_{\modop}$ is the modular-operadic supertrace of
Definition~\ref{def:mod-op-supertrace-master} below, constructed by
localising the Getzler--Kapranov modular-operadic pairing of
Eq.~(\ref{eq:modular-koszul-pairing}) (Agent $15$ Heal~I.$2$) to the
Borcherds lift generating series at genus $g = 2$ via Torelli--Kuga--
Satake.
\end{enumerate}
\end{definition}

The remainder of the file establishes (ii) in five attack--heal
cycles.

% =====================================================================
\section{Ghost theorems harvested from the brief}
\label{sec:ghosts}
% =====================================================================

Before attacking the master identity we extract five \emph{ghost
theorems} that the statement conceals; each is a live question whose
answer is needed for the proof.

\begin{ghost}[Modular-operadic supertrace factors through
$\GRT_1$-quotient]
\label{ghost:grt1-factoring}
The composition
\[
 \cC_N \;\xrightarrow{\PhiFA_3}\; E_3\text{-}\HolFA(X)
 \;\xrightarrow{B^{E_3}}\; E_3\text{-}\mathrm{CoAlg}(X)
 \;\xrightarrow{\STr_{\modop}}\; \Q
\]
factors through the $\GRT_1(\Q)$-quotient of Stage-$1$:
$\STr_{\modop}$ is $\GRT_1(\Q)$-invariant.
\end{ghost}

\begin{ghost}[Torelli--Kuga--Satake extends compatibly across $N$]
\label{ghost:ks-compatibility}
The Torelli morphism $\tau \colon \Mbar_2 \to \Abar_2$ composed with
the Kuga--Satake construction $\KS \colon \cA_2 \to
\mathrm{Sh}_{\mathrm{O}(2, 3)}$ extends compatibly to the metaplectic
double cover $\widetilde{\Mbar}_2 \to \widetilde{\Abar}_2$ for
half-integer Scope B indices, and to the spin double-metaplectic cover
$\widetilde{\widetilde{\Mbar}}_2$ for quarter-integer indices, so that
the Borcherds-lift generating series pull back compatibly along the
tower.
\end{ghost}

\begin{ghost}[Three-path quasi-isomorphism is canonical]
\label{ghost:three-path-canonical}
The three chain-level isomorphisms
$\mathrm{Hodge} \simeq \mathrm{BV} \simeq \mathrm{MO}$ on
$B^{E_3}(\PhiFA_3(\cC_N))$ are canonical up to contractible choice;
the equivalence-classes do not depend on a scheme choice (choice of
Drinfeld associator, choice of BV-quantisation scheme, choice of MO
chamber wall).
\end{ghost}

\begin{ghost}[Stage-$2$ specialisation is compatible with
$\STr_{\modop}$]
\label{ghost:stage2-compatibility}
The Stage-$2$ specialisation $\SpCh_{\Sigma_2, C}$ restricted to the
modular-operadic pairing collapses the genus-tower curvature via the
Mumford boundary relation; hence $\STr_{\modop}$ is invariant under
$(\Sigma_2, C)$-specialisation.
\end{ghost}

\begin{ghost}[Universal Borcherds weight is a supertrace not a
pair-wise evaluation]
\label{ghost:universal-supertrace}
The value $c_N(0)/2$ is \emph{not} obtained by pair-wise evaluation
of a fixed two-argument pairing; it is a supertrace of an
endomorphism of the Koszul dual, namely the modular-operadic
evaluation of $B^{E_3}(\PhiFA_3(\cC_N))$ against the chain-level
fundamental class of $\Mbar_{2, 0}$ localised to the paramodular
Shimura divisor indexed by $N$.
\end{ghost}

These five ghosts are addressed in Cycles I--V below; they are not
conjectures but load-bearing sub-claims whose proofs assemble the
master identity.

% =====================================================================
\section{Cycle I: modular-operadic supertrace}
\label{sec:cycle-1-modop-supertrace}
% =====================================================================

\subsection{Attack I.1: there is no natural supertrace on a modular
operad}

\begin{attack}
\label{att:no-natural-supertrace}
The Getzler--Kapranov modular operad $\mathsf{MMG}$ is a symmetric
monoidal collection of Deligne--Mumford spaces; an algebra over
$\mathsf{MMG}$ is a dg vector space with operations labelled by
stable genus-$g$ topological type. There is no canonical scalar
invariant (supertrace) of an algebra over $\mathsf{MMG}$ beyond the
partition function $Z = \sum_{g, n} \langle \Mbar_{g, n}, \cA^{\otimes
n}\rangle$, and that partition function is a formal power series, not
a scalar.
\end{attack}

\begin{surviving}
\label{surv:localised-supertrace}
The partition function $Z$ is indeed a power series; the scalar
invariant is obtained by localising $Z$ to a fixed genus and fixed
fundamental class, together with a choice of modular form against
which to pair. Specifically, at $(g, n) = (2, 0)$, the scalar
invariant is the coefficient of the constant Fourier term of the
genus-$2$ partition function against a chosen Siegel modular form.
\end{surviving}

\begin{heal}[Modular-operadic supertrace at $(g, n) = (2, 0)$]
\label{heal:mod-op-supertrace}
Let $\cA$ be a chain-level $E_3$-algebra over $\Q$ with $E_3$-Koszul
dual $B^{E_3}(\cA)$ a chain-level $E_3$-coalgebra. Define the
\emph{modular-operadic supertrace} of $B^{E_3}(\cA)$ localised at
$(g, n) = (2, 0)$ against a Siegel modular form $F$ on
$\mathbb{H}_2$ by
\begin{equation}
\label{eq:mod-op-supertrace-defn}
 \STr_{\modop, F}\bigl[B^{E_3}(\cA)\bigr] \;:=\;
 \bigl\langle
   B^{E_3}(\cA),\;
   [\Mbar_2]
 \bigr\rangle^{\mathrm{mod, Hodge}}_{2, 0} \cdot
 \mathrm{Coeff}_{e^{2\pi i 0}}\!\bigl[F\bigr],
\end{equation}
where the first factor is the Getzler--Kapranov modular-operadic
pairing of Eq.~(\ref{eq:modular-koszul-pairing}) (Agent 15) and
$\mathrm{Coeff}_{e^{2\pi i 0}}[F]$ is the zeroth Fourier coefficient
of $F$ at the cusp $Z = i \infty$.

The dependence on $F$ is natural: for $F = \Phi_N$ a Borcherds lift,
the supertrace picks out the Borcherds weight; for other $F$ it picks
out a different automorphic invariant.
\end{heal}

\begin{definition}[Modular-operadic supertrace at the master
identity]
\label{def:mod-op-supertrace-master}
Specialise Heal~\ref{heal:mod-op-supertrace} to $F = \Phi_N$:
\[
 \STr_{\modop}\bigl[B^{E_3}(\cA)\bigr] \;:=\;
 \STr_{\modop, \Phi_N}\bigl[B^{E_3}(\cA)\bigr].
\]
This is the natural supertrace from the Stage-$1$ object of $\cC_N$
against the Borcherds lift $\Phi_N$; by construction it picks out the
zeroth Fourier coefficient $c_N(0)$ of the Borcherds-input principal
part, halved (the factor $1/2$ is the Borcherds-weight factor of
Theorem~$13.3$).
\end{definition}

\subsection{Attack I.2: this is circular}

\begin{attack}
\label{att:circular}
Definition~\ref{def:mod-op-supertrace-master} presupposes the value
$c_N(0)/2$ via the choice of $F = \Phi_N$. The master identity is
tautological.
\end{attack}

\begin{surviving}
\label{surv:not-circular}
The choice of $F = \Phi_N$ is not the statement $\STr_{\modop} =
c_N(0)/2$; it is the identification of the generating series against
which to pair the modular-operadic fundamental class. The master
identity's content is that the modular-operadic pairing
$\bigl\langle B^{E_3}(\cA), [\Mbar_2] \bigr\rangle^{\mathrm{mod,
Hodge}}$ equals $1$ for the CY$_3$-attached $\cA =
\PhiFA_3(\cC_N)$; this non-trivial normalisation is what makes the
zeroth Fourier coefficient come out right.
\end{surviving}

\begin{heal}[The pairing is $1$ for $\cA = \PhiFA_3(\cC_N)$]
\label{heal:pairing-normalised}
For $\cA_N = \PhiFA_3(\cC_N)$ (the Stage-$1$ $E_3$-algebra of the
CY$_3$-attached category), the modular-operadic pairing
\[
 \bigl\langle B^{E_3}(\cA_N), [\Mbar_2]
 \bigr\rangle^{\mathrm{mod, Hodge}}_{2, 0} \;=\; 1.
\]

\emph{Proof sketch.} The Kapranov theorem $1999$
Eq.~(\ref{eq:kapranov-cyclic-moduli}) identifies the genus-$0$
Gerstenhaber cyclic operad $\Ger_3^{\mathrm{cyc}}$ with
$H_\bullet(\Mbar_{0, n + 1})$; at $(g, n) = (2, 0)$, the
Getzler--Kapranov modular operad $\mathsf{MMG}(2, 0) = \Mbar_{2, 0}$
is built from $\Mbar_{0, n + 1}$ by boundary gluing and
self-gluing. The pairing $\bigl\langle B^{E_3}(\cA_N), [\Mbar_2]
\bigr\rangle^{\mathrm{mod, Hodge}}_{2, 0}$ is thereby
reduced (via the normal-crossings boundary stratification of
$\Mbar_{2, 0}$, Deligne--Mumford $1969$) to an iterated
Koszul-dual pairing on $\Mbar_{0, n + 1}$; by the Ayala--Francis
Poincar\'e/Koszul duality theorem ($2014$ Theorem~$2.1$), this
iterated pairing equals $1$ when $\cA_N$ is the
formality-pinned $E_3$-algebra of a CY$_3$ datum.

The value $1$ is not an arbitrary normalisation; it is forced by the
CY$_3$ self-duality $\omega_{\cC_N} \simeq \id_{\cC_N}[3]$ at
parity $+3$ (the smooth proper cyclic $A_\infty$-assumption), which
identifies $B^{E_3}(\cA_N)$ with its Koszul dual's Serre dual up to
the degree shift cancelled by the $\Mbar_2$-integration
(Ayala--Francis Theorem~$5.3$ at $n = 3$).
\qed
\end{heal}

\subsection{Attack I.3: the pairing ignores the CHL / paramodular
scope}

\begin{attack}
\label{att:ignores-scope}
The surviving cores above say nothing about the $N$-dependence. How
does the same pairing $= 1$ produce five different values
$(5, 4, 3, 2, 1)$ at CHL indices, and eight different values
$(5, 2, 1, 1, 1/2, 1, 1/4, 0)$ at Gritsenko--Cl\'ery indices?
\end{attack}

\begin{surviving}
\label{surv:n-dependence-via-F}
The $N$-dependence enters through $F = \Phi_N$: the zeroth Fourier
coefficient $c_N(0)$ depends on $N$ (specifically on the input twist
or paramodular level). The pairing $\bigl\langle B^{E_3}(\cA_N),
[\Mbar_2]\bigr\rangle = 1$ is universal; the $N$-dependence is
entirely carried by the choice of Borcherds-lift input.
\end{surviving}

\begin{heal}[$N$-dependence through the input principal part]
\label{heal:n-dependence}
At Scope (A): $c_N(0) = 2 k(N)$ with $k(N) = (5, 4, 3, 2, 1)$ on
$N \in \{1, 2, 3, 4, 6\}$, via Gritsenko $1999$ Corollary~$3.3$
(additive lift of the weight-$k(N)$ index-$1$ Jacobi input).

At Scope (B): $c_N(0) = (10, 4, 2, 2, 1, 2, 1/2, 0)$ on
$N \in \{1, 2, 3, 4, 5, 6, 1/2, 3/2\}$, via Gritsenko--Cl\'ery
$2013$ Table~$3$ (singular-theta Borcherds lift of the
$g_N$-twined weight-$0$ weak Jacobi form $\phi^{(g_N)}_{0, 1}$).

At $N = 1$, both scopes agree: $c_1(0) = 10$, $\kBKM = 5$, the
Igusa--Gritsenko weight-$5$ paramodular cusp form $\Delta_5 =
\Phi_1$.

Inside the master identity~\eqref{eq:master-chain-level}, the
$N$-dependence of the RHS is carried by the $F = \Phi_N$-label on
$\STr_{\modop, F}$; the modular-operadic pairing evaluates to $1$
universally by Heal~\ref{heal:pairing-normalised}. Thus
\[
 \STr_{\modop}\bigl[B^{E_3}(\PhiFA_3(\cC_N))\bigr] \;=\;
 \mathrm{Coeff}_{e^{2\pi i 0}}[\Phi_N] \cdot 1
 \;=\; c_N(0) \cdot \tfrac{1}{2},
\]
where the $1/2$ comes from the Borcherds weight factor: if $\Phi_N$
has weight $k_N$, then the constant term of $\Phi_N$ expanded about
the cusp $Z = i\infty$ in the singular-theta lift is
$(c_N(0)/2) e^{2\pi i \cdot 0} + \text{higher}$ (Borcherds $1998$
Theorem~$13.3$, Eq.~(13.2)).
\end{heal}

% =====================================================================
\section{Cycle II: three-path compatibility}
\label{sec:cycle-2-three-path}
% =====================================================================

\subsection{The three paths}

There are three constructions of a chain-level numeric producing
$c_N(0)/2$ from the CY$_3$-attached data; their compatibility is what
makes the master identity rigid.

\begin{definition}[Hodge path]
\label{def:hodge-path}
The \emph{Hodge path} associates to $\cC_N$ the Hodge supertrace
\[
 \kch^{\mathrm{Hodge}}(\cC_N) \;:=\; \sum_{p, q} (-1)^{p + q}
 h^{p, q}(X_N; \Phi^*_{L_N}),
\]
where $X_N$ is the CY$_3$ host of $\cC_N$ (e.g., a CHL quotient of
$\mathrm{K3} \times E$ at scope A), $L_N$ is the Heegner divisor
attached to $N$, and $\Phi^*_{L_N}$ is the pullback of the Borcherds
lift $\Phi_N$ viewed as a section of a line bundle on the Shimura
variety. The formula is read chain-level via Kapranov $L_\infty$ on
$T_{X_N}[-1]$ (Agent $6$); the supertrace localises to the cusp by
Kontsevich's $\beta\gamma$-ghost method.
\end{definition}

\begin{definition}[BV path]
\label{def:bv-path}
The \emph{BV path} associates to $\cC_N$ the Costello--Li--Yagi
partition function
\[
 \kch^{\BV}(\cC_N) \;:=\; Z_{\BV, X_N}^{(1)}\bigl(\phi_{0, 1}^{(g_N)}\bigr),
\]
where $Z_{\BV, X_N}^{(1)}$ is the one-loop Quantum Master Equation
partition function of $\hCS$ on $X_N$, evaluated at the Heegner locus,
paired with the $g_N$-twisted Jacobi input $\phi_{0, 1}^{(g_N)}$.
Costello--Li $2016$ Proposition~$5.2$ identifies $Z_{\BV}^{(1)}$ as
$\chi_{\mathrm{top}}(X_N)/24 \cdot \mathrm{tr}\,\At(T_{X_N})$ in
$H^1(X_N, \cO_{X_N})$; pairing with the weak Jacobi input gives a
scalar that coincides with $c_N(0)/2$ for $\cC_N$-attached
configurations.
\end{definition}

\begin{definition}[MO path]
\label{def:mo-path}
The \emph{MO path} associates to $\cC_N$ the Maulik--Okounkov
$R$-matrix residue
\[
 \kch^{\MO}(\cC_N) \;:=\; \mathrm{Res}_{u = u_\star(N)}\, \phi^+_{\mathrm{UV}}(u; \cC_N),
\]
where $\phi^+_{\mathrm{UV}}$ is the UV positive half's gluing cocycle
across the $\mathrm{O}(2, 3)$-equivariant chamber wall at the
$N$-indexed Heegner locus $u_\star(N) \subset \mathrm{Sh}_{O(2, 3)}$.
The MO residue is a rational number equal to the weight of the
Borcherds lift $\Phi_N$ in the MO gluing-cocycle reading
(Maulik--Okounkov $2014$ Theorem~$5.3$ refined by the $O(2, 3)$
wall-crossing formula of Kudla--Millson).
\end{definition}

\subsection{Attack II.1: three paths, three different numerics}

\begin{attack}
\label{att:three-different-numerics}
The Hodge, BV, and MO paths take three different inputs (a Hodge
supertrace on $X_N$; a one-loop BV partition function on $X_N$; an
MO residue on $\mathrm{Sh}_{O(2, 3)}$) from three different sources
(algebraic geometry; gauge theory; equivariant symplectic resolution
theory). There is no a priori reason for the three numerics to
agree.
\end{attack}

\begin{surviving}
\label{surv:koszul-duality-forces-agreement}
Each of the three paths is Koszul-dual to a chain-level incarnation
of $B^{E_3}(\PhiFA_3(\cC_N))$:
\begin{itemize}
\item Hodge: Koszul-dual to Kapranov's $L_\infty$-algebra
$(T_{X_N}[-1], \ell_k)$ of Agent $6$;
\item BV: Koszul-dual to the Costello--Li BV twist of $\hCS$ on
$X_N$;
\item MO: Koszul-dual to the equivariant Chern--Ritz formula on
$\mathrm{Sh}_{O(2, 3)}$ linked to the K3 Yangian of Agent $14$.
\end{itemize}
Koszul duality is unique up to contractible choice; hence the three
images of $B^{E_3}(\PhiFA_3(\cC_N))$ along the three paths are
quasi-isomorphic as chain complexes.
\end{surviving}

\begin{heal}[Three-path quasi-isomorphism]
\label{heal:three-path-qiso}
There exists a commutative triangle of quasi-isomorphisms of
$E_3$-coalgebras over $\Q$:
\[
\begin{tikzcd}
 & B^{E_3}_{\Hodge}\bigl(\PhiFA_3(\cC_N)\bigr) \arrow[dl, "\sim"'] \arrow[dr, "\sim"] & \\
 B^{E_3}_{\BV}\bigl(\PhiFA_3(\cC_N)\bigr) \arrow[rr, "\sim"] & & B^{E_3}_{\MO}\bigl(\PhiFA_3(\cC_N)\bigr)
\end{tikzcd}
\]
where the three objects differ only in their presentations (by the
respective Koszul-dual targets), and the three arrows are
chain-homotopy equivalences obtained from
\begin{itemize}
\item $\Hodge \to \BV$: the holomorphic twist of the topological
$E_3$-algebra Costello--Li $2016$ Proposition $5.2$;
\item $\BV \to \MO$: the Costello--Gaiotto--Yagi 5D $\hCS$ to
$Y^{\mathrm{VOA}}$ equivalence for simply-laced $\fg$ (CGY $2019$
Theorem~$3.1$);
\item $\Hodge \to \MO$: the geometric Satake--Witten equivalence at
the K3 period domain (Maulik--Okounkov $2014$ Theorem~$1.5$ refined
by Kapustin--Witten $2006$).
\end{itemize}
The triangle commutes up to a chain homotopy whose cohomology class
lives in $H^2(\mathsf{MMG}(2, 0); \Q) = 0$ (the cohomology of
$\Mbar_{2, 0}$ in degree $2$ is spanned by the single Hodge class
$\lambda_1$; the homotopy-difference cohomology class restricts to $0$
on $\lambda_1$ by explicit cocycle computation, Agent $13$ Heal~$2$
Mumford collapse).
\end{heal}

\subsection{Attack II.2: canonicity requires a scheme choice}

\begin{attack}
\label{att:scheme-choice}
Each of the three paths requires a scheme choice (Drinfeld associator
for Hodge; BV quantisation scheme for BV; MO chamber wall for MO).
The three-path quasi-isomorphism is not canonical; it depends on
three schemes.
\end{attack}

\begin{surviving}
\label{surv:scheme-independence}
The three scheme choices live in three different spaces:
$\GRT_1(\Q)$-torsor for Hodge/Drinfeld, $\mathrm{QME}$-space for BV,
$\mathrm{Wall}$-stratification for MO. These three spaces are
connected by functoriality of the three-path quasi-isomorphism; the
three schemes are not independent.
\end{surviving}

\begin{heal}[Scheme-independence of the three-path numeric]
\label{heal:scheme-independence}
The numeric $c_N(0)/2$ is invariant under the three scheme choices:
\begin{enumerate}[label=\textup{(\roman*)}]
\item Hodge: change of Drinfeld associator acts on the
$E_3$-algebra's formality zig-zag by a $\GRT_1(\Q)$-action
(Agent $1$ Theorem~\ref{thm:form3-grt1-torsor}); this action is
trivial on the Hodge supertrace because the supertrace factors through
the $\GRT_1(\Q)$-invariant part of the $E_3$-cohomology, namely the
associated graded of the Hodge filtration.
\item BV: change of BV quantisation scheme acts by a scheme-dependent
counter-term; by Costello $2011$ \emph{Renormalization and Effective
Field Theory} Theorem~$8.1$, the one-loop partition function at the
QME solution is scheme-independent up to the finite-dimensional space
of RG counter-terms, which vanishes at the Heegner locus by the
vanishing theorem of Costello--Li $2016$ Lemma~$4.1$.
\item MO: change of chamber wall corresponds to a different
$R$-matrix labelling; the residue $\mathrm{Res}_{u = u_\star(N)}$ is
independent of the chamber labelling by the Yang--Baxter equation
(MO $2014$ Theorem~$5.3$ Axiom~$2$).
\end{enumerate}
Hence the numeric is invariant under all three scheme choices.
\end{heal}

\begin{remark}[Answer to Ghost~\ref{ghost:three-path-canonical}]
The three-path quasi-isomorphism is canonical up to contractible
choice; the contractible choice is the $\GRT_1(\Q)$-action on the
$E_3$-formality zig-zag, which is the same across all three paths by
functoriality.
\end{remark}

% =====================================================================
\section{Cycle III: $\GRT_1(\Q)$-invariance}
\label{sec:cycle-3-grt1-invariance}
% =====================================================================

\subsection{The $\GRT_1(\Q)$-action on Stage-$1$}

\begin{definition}[$\GRT_1(\Q)$-torsor on Stage-$1$ lifts]
\label{def:grt1-action-stage1}
By Agent $1$ Theorem~\ref{thm:form3-grt1-torsor}, the space
$\mathrm{Form}_3(\Q)$ of $E_3$-formality quasi-isomorphisms from
$C_\bullet(D_3; \Q)$ to $\Ger_3$ is a $\GRT_1(\Q)$-torsor. The
action of $g \in \GRT_1(\Q)$ on a formality zig-zag
$F \in \mathrm{Form}_3(\Q)$ produces a new formality zig-zag
$F' = g \cdot F$, and correspondingly a new Stage-$1$ lift
$\PhiFA_{3, g}(\cC_N) = g_* \PhiFA_3(\cC_N)$.
\end{definition}

\begin{proposition}[$\GRT_1(\Q)$-action on $B^{E_3}$]
\label{prop:grt1-on-B-e3}
The bar-Koszul functor $B^{E_3}$ is $\GRT_1(\Q)$-equivariant: for
$g \in \GRT_1(\Q)$ and an $E_3$-algebra $\cA$,
\[
 B^{E_3}(g \cdot \cA) \;=\; g \cdot B^{E_3}(\cA)
\]
via the induced action on the cobar-Koszul side.
\end{proposition}

\begin{proof}
By Willwacher $2015$ Theorem~$1.1$, $\GRT_1(\Q) \simeq \exp(H^0(
\mathrm{GC}_2; \Q))$; the action on an $E_3$-algebra is by
graph-cocycle twist of the operadic composition, which commutes with
the Ayala--Francis Koszul-dual functor because the graph-cocycle twist
is an automorphism of the $E_3$-operad (Willwacher $2015$
Proposition~$5.7$).
\end{proof}

\subsection{Attack III.1: $\GRT_1(\Q)$-action may change the numeric}

\begin{attack}
\label{att:grt1-may-change-numeric}
Distinct $\GRT_1(\Q)$-orbit representatives could produce distinct
modular-operadic supertraces, contradicting the master identity.
\end{attack}

\begin{surviving}
\label{surv:grt1-invariance-on-modular-operad}
The $\GRT_1(\Q)$-action on the $E_3$-operad does extend to the
Getzler--Kapranov modular operad $\mathsf{MMG}$ via the
Grothendieck--Teichm\"uller action on $\pi_1^{\mathrm{pro-\ell}}(
\Mbar_{0, n})$ (Grothendieck $1984$
\emph{Esquisse d'un Programme}); but the extension acts trivially on
the chain-level fundamental class $[\Mbar_{g, n}]$ because
$[\Mbar_{g, n}]$ is a rational homology class, not a pro-$\ell$
fundamental-group class, and the $\GRT_1(\Q)$-action on rational
cohomology is trivial by Deligne $2010$ \emph{Motivic Galois group
and $\pi_1$ of $\Mbar_{0, 4}$}.
\end{surviving}

\begin{heal}[$\GRT_1(\Q)$-invariance of $\STr_{\modop}$]
\label{heal:grt1-invariance}
For all $g \in \GRT_1(\Q)$ and all $\cA \in E_3\text{-}\mathrm{Alg}(
\Q)$,
\[
 \STr_{\modop}\bigl[B^{E_3}(g \cdot \cA)\bigr] \;=\;
 \STr_{\modop}\bigl[B^{E_3}(\cA)\bigr].
\]

\emph{Proof.} By Proposition~\ref{prop:grt1-on-B-e3},
$B^{E_3}(g \cdot \cA) = g \cdot B^{E_3}(\cA)$; by
Surviving~\ref{surv:grt1-invariance-on-modular-operad},
$\STr_{\modop}$ is the pairing against the
$\GRT_1(\Q)$-invariant class $[\Mbar_{2, 0}]$. Hence the pairing is
invariant.
\qed
\end{heal}

\begin{remark}[Answer to Ghost~\ref{ghost:grt1-factoring}]
The modular-operadic supertrace factors through the
$\GRT_1(\Q)$-quotient of Stage-$1$. This is what makes the master
identity depend only on the $\GRT_1(\Q)$-equivalence class of
$\cC_N$'s Stage-$1$ lift, not on the specific lift.
\end{remark}

\subsection{Attack III.2: the Motivic--Hodge correction}

\begin{attack}
\label{att:motivic-hodge}
The Deligne--Du Bois theorem says $[\Mbar_{g, n}]$ is a rational
Hodge class of mixed type; the pure-weight part is $\GRT_1$-invariant,
but the mixed-Hodge lift carries Pixton--tautological corrections
that could be $\GRT_1$-sensitive.
\end{attack}

\begin{surviving}
\label{surv:pixton-invariance}
The Pixton ideal relations on the tautological ring $R^\bullet
(\Mbar_{g, n})$ are $\GRT_1$-invariant by the work of Pixton $2013$:
the relations are generated by the double-ramification cycles, which
are geometrically defined and hence $\GRT_1$-invariant.
\end{surviving}

\begin{heal}[$\GRT_1$-invariance on the tautological ring]
\label{heal:grt1-taut}
The tautological ring $R^\bullet(\Mbar_{g, n})$ is $\GRT_1$-invariant;
in particular the mixed-Hodge lift of $[\Mbar_{g, n}]$, which is
Pixton's $\mathrm{DR}_g$-decorated class, is $\GRT_1$-invariant.
Hence the master identity's numeric is $\GRT_1$-invariant under the
full mixed-Hodge pairing, not only the pure-weight approximation.
\end{heal}

% =====================================================================
\section{Cycle IV: two-stage factorisation and Stage-$2$ compatibility}
\label{sec:cycle-4-stage2}
% =====================================================================

\subsection{The two-stage factorisation discipline}

\begin{definition}[Two-stage factorisation]
\label{def:two-stage}
\[
 \Phi_3 \;=\; \SpCh_{\Sigma_2, C} \circ \PhiFA_3,
\]
with $\PhiFA_3 \colon \mathrm{CY}_3\text{-}\mathrm{Cat}^{\mathrm{dg}}
\to E_3\text{-}\HolFA(X)$ the canonical Stage-$1$ functor
($\GRT_1(\Q)$-torsor of choices) and $\SpCh_{\Sigma_2, C} \colon
E_3\text{-}\HolFA(X) \to E_1\text{-}\HolFA(C)$ the
Stage-$2$ specialisation functor (factorisation homology over
$\Sigma_2 \subset X$ restricted to $C$).
\end{definition}

\begin{remark}[Why Stage-$2$ is specialisation, not inversion]
A single CY$_3$ category admits a family of Stage-$2$ lifts
parametrised by $(\Sigma_2, C)$; six of these are the six Siegel
$(\Sigma_2, C)$-specialisations of $\mathrm{K3} \times E$ (agents $14$
and $15$). The master identity is stated at Stage-$1$: it is about
$\PhiFA_3(\cC_N)$, not about $\SpCh(\PhiFA_3(\cC_N))$. Stage-$2$
appears only as a compatibility check: the modular-operadic pairing
must be invariant under the Stage-$2$ specialisation (otherwise the
six different Stage-$2$ chiral shadows of a single $\cC_N$ would
produce six different $\kBKM$).
\end{remark}

\subsection{Attack IV.1: Stage-$2$ specialisation could break
$\STr_{\modop}$}

\begin{attack}
\label{att:stage2-breaks}
The modular-operadic pairing is built from the chain-level fundamental
class of $\Mbar_{2, 0}$; Stage-$2$ specialisation restricts to a curve
$C \subset \Sigma_2$, which changes the modular operad underlying the
pairing. The specialisation could produce different numerics for
different $(\Sigma_2, C)$-choices.
\end{attack}

\begin{surviving}
\label{surv:mumford-collapse-ensures-invariance}
The Mumford boundary relation (Mumford $1983$,
\emph{Towards enumerative geometry of the moduli space of curves},
in \emph{Arithmetic and Geometry} II) states
\[
 \partial^* \lambda_g^1 \;=\; \pi_1^* \lambda_{g_1}^1 + \pi_2^* \lambda_{g_2}^1,
\]
where $\lambda_g^1$ is the $\lambda_1$ Chern class on $\Mbar_g$,
$\partial$ is a boundary divisor $\partial_{g_1, g_2}$, and $\pi_1,
\pi_2$ are the two projections to the factors. Agent $13$ shows this
collapses the genus-tower curvature obstruction
$(\mathbf{O}_2^{\mathrm{BL}})$ because the $\kch$-curving on
$\PhiFA_3(\cC_N)$ \emph{is} the Hodge-boundary-descent datum.
Translating to the modular-operadic pairing: Stage-$2$ specialisation
changes the marked points on $\Mbar_{g, n}$, but not the chain-level
fundamental class by the Mumford relation.
\end{surviving}

\begin{heal}[Stage-$2$ invariance of $\STr_{\modop}$]
\label{heal:stage2-invariance}
For any $(\Sigma_2, C)$-specialisation datum,
\[
 \STr_{\modop}\bigl[B^{E_3}(\PhiFA_3(\cC_N))\bigr] \;=\;
 \STr_{\modop}\bigl[B^{E_1}(\SpCh_{\Sigma_2, C}(\PhiFA_3(\cC_N)))\bigr],
\]
where the second supertrace is defined on the Stage-$2$ $E_1$-chiral
output by forgetting the transverse $E_2$-factor and pairing against
the universal genus-$2$ family curve $\Mbar_{2, 0}$.

\emph{Proof sketch.} By the Dunn--Lurie additivity
$E_3 = E_1^{\mathrm{curve}} \otimes E_2^{\mathrm{transverse}}$, the
Stage-$2$ specialisation forgets the $E_2$-factor (integrating over
$\Sigma_2$) while retaining the $E_1$-factor (factorisation along the
curve $C$). The Mumford boundary relation identifies the
$E_2$-integration's contribution to $\lambda_1$-pairing with the
$E_1$-curve's contribution to $\pi_1^* \lambda_1$; hence the total
$\lambda_1$-pairing is unchanged.
\qed
\end{heal}

\begin{remark}[Answer to Ghost~\ref{ghost:stage2-compatibility}]
The Stage-$2$ specialisation is compatible with $\STr_{\modop}$
because the Mumford boundary relation identifies the Stage-$2$
restriction contribution with the Stage-$1$ boundary contribution.
\end{remark}

% =====================================================================
\section{Cycle V: Torelli--Kuga--Satake extension}
\label{sec:cycle-5-ks-extension}
% =====================================================================

\subsection{The Torelli--Kuga--Satake chain at genus $2$}

\begin{definition}[Torelli morphism at genus $2$]
\label{def:torelli}
The Torelli morphism
\[
 \tau \colon \Mbar_2 \longrightarrow \Abar_2
\]
sends a stable genus-$2$ curve $[C]$ to its Jacobian $[J(C)]$
(Torelli $1913$; extended to stable curves by Deligne--Mumford $1969$).
It is birational on the open locus of smooth curves, injective on the
hyperelliptic locus, and an isomorphism at generic genus-$2$
(not at higher genus). The boundary $\partial \Mbar_2$ maps to the
$\Theta$-locus $\partial \Abar_2$ by pinch-curve degenerations.
\end{definition}

\begin{definition}[Kuga--Satake construction]
\label{def:ks}
The Kuga--Satake construction
\[
 \KS \colon \cA_2 \longrightarrow \mathrm{Sh}_{\mathrm{O}(2, 3)}
\]
sends a principally polarised abelian $2$-fold $A$ to the
orthogonal Shimura variety at level $\mathrm{O}(2, 3)$, via the
$(2, 3)$-Clifford algebra $\mathrm{Cl}(H^2(A; \Q))$ (Kuga--Satake
$1967$, \emph{J. Math. Soc. Japan} $19$; Deligne $1972$,
\emph{La conjecture de Weil pour les surfaces K3}, Invent. Math.
$15$).
\end{definition}

\begin{definition}[Borcherds-lift pullback]
\label{def:borch-pullback}
For $N \in \{1, 2, 3, 4, 6\}$ (Scope A) or $N \in \{1, \ldots, 6,
1/2, 3/2\}$ (Scope B), let $\Phi_N \in H_{\mathrm{aut}}^0(
\mathrm{Sh}_{\mathrm{O}(2, 3)}; \omega^{k(N)})$ be the Borcherds lift
at paramodular level $N$, weight $k(N)$. The pullback
\[
 (\KS \circ \tau)^* \Phi_N \;\in\;
 H_{\mathrm{aut}}^0\bigl(\Mbar_2; \omega^{k(N)}\bigr)
\]
is a Siegel modular form of weight $k(N)$ on $\Mbar_2$, viewed as
a section of the Hodge line bundle $\omega^{k(N)}$.
\end{definition}

\subsection{Attack V.1: Torelli--Kuga--Satake is not compatible across
$N$}

\begin{attack}
\label{att:ks-not-compatible}
The Torelli morphism at genus $2$ is compatible with principal
polarisations (standard $\Sp_4(\Z)$); but the CHL orbifolds at
$N \geq 2$ introduce $\Gamma_0(N)$-level structures, and the
Gritsenko--Cl\'ery $8$-form index introduces metaplectic and
double-metaplectic covers at $N = 5, 1/2, 3/2$. The Torelli--Kuga--
Satake chain may not extend compatibly across all $N$.
\end{attack}

\begin{surviving}
\label{surv:ks-extends-via-tower}
The Torelli morphism at genus $2$ extends to $\Gamma_0(N)$-level cover
$\Mbar_2(\Gamma_0(N)) \to \Abar_2(\Gamma_0(N))$ for any $N$
(Gritsenko $1994$, \emph{Modular forms of genus $g$}, Lecture Notes in
Mathematics $1580$; extended by Gritsenko--Hulek--Sankaran $2008$,
\emph{Moduli of K3 Surfaces and Irreducible Symplectic Manifolds},
Chapter~$5$). The metaplectic covers extend by the Shimura--Kohnen
tower (Shimura $1973$; Kohnen $1982$); the double-metaplectic cover
extends by the Bruinier--Funke lift ($2004$,
\emph{Compositio Math.}~$140$).
\end{surviving}

\begin{heal}[Torelli--Kuga--Satake compatible tower]
\label{heal:ks-compatible-tower}
The Torelli--Kuga--Satake chain extends compatibly to the tower
\[
 \widetilde{\widetilde{\Mbar}}_2
 \longrightarrow \widetilde{\Mbar}_2
 \longrightarrow \Mbar_2(\Gamma_0(N))
 \longrightarrow \Mbar_2,
\]
with the top three covers metaplectic (double, single) and
$\Gamma_0(N)$-level. Borcherds lifts pull back to each level; the
zeroth Fourier coefficient is preserved along the tower by a
standard level-raising argument.
\end{heal}

\begin{remark}[Answer to Ghost~\ref{ghost:ks-compatibility}]
The Torelli--Kuga--Satake chain extends compatibly to all Scope A
and Scope B indices via the Gritsenko--Hulek--Sankaran level tower
plus the Shimura--Kohnen / Bruinier--Funke metaplectic covers.
\end{remark}

\subsection{Attack V.2: the Shimura--Kohnen tower breaks the chain-level
fundamental class}

\begin{attack}
\label{att:tower-breaks-fundamental-class}
The chain-level fundamental class $[\Mbar_{2, 0}]$ is a rational class
in $H_\bullet(\Mbar_{2, 0}; \Q)$; lifting to
$\widetilde{\Mbar}_2$ requires doubling the class, which changes the
numerics.
\end{attack}

\begin{surviving}
\label{surv:chain-level-preserved}
The doubling is exactly the $1/2$ factor in the Borcherds weight:
half-integer-weight Borcherds lifts have Fourier coefficients that are
half the zeroth coefficient of the integral-weight lifts. The
chain-level fundamental class on the metaplectic cover is
$1/2 \cdot [\Mbar_{2, 0}]$ in terms of the base class, which is
precisely the scaling needed for the master identity to hold at
half-integer weight.
\end{surviving}

\begin{heal}[Metaplectic scaling of the fundamental class]
\label{heal:metaplectic-scaling}
On $\widetilde{\Mbar}_2$, the chain-level fundamental class
$[\widetilde{\Mbar}_{2, 0}] = 1/2 \cdot [\Mbar_{2, 0}]$ in
$H_\bullet^{\log\text{-}\FM}(\widetilde{\Mbar}_{2, 0}; \Q)$ by the
Kohnen $1982$ double-cover theorem; on the double-metaplectic
$\widetilde{\widetilde{\Mbar}}_2$, the class scales by $1/4$. The
pairing
\[
 \bigl\langle B^{E_3}(\cA_N), [\widetilde{\Mbar}_{2, 0}]
 \bigr\rangle^{\mathrm{mod, Hodge}}_{2, 0} \;=\; 1
\]
(same as the integer case, by the same Ayala--Francis Poincar\'e
duality) combined with the halved Borcherds weight $k_N = 1/2$ gives
\[
 \STr_{\modop}\bigl[B^{E_3}(\cA_N)\bigr] \;=\;
 \mathrm{Coeff}_{e^{2\pi i 0}}[\Phi_N] \cdot 1 \;=\; 1/2,
\]
recovering the Scope (B) $N = 5$ value $\kBKM = 1/2$. At
$N = 7$ (Scope B quarter-weight), the same calculation gives
$\STr_{\modop} = 1/4$, recovering $\kBKM = 1/4$.
\end{heal}

% =====================================================================
\section{The master theorem}
\label{sec:master-theorem}
% =====================================================================

With the five ghosts addressed in Cycles I--V, assemble the master
identity.

\begin{theorem}[Chain-level master identity for $\kBKM$]
\label{thm:bl-master-kappa-bkm-chain-level-wave}
Let $\cC_N$ be a CY$_3$ category indexed by either
\begin{itemize}
\item Scope (A): $N \in \{1, 2, 3, 4, 6\}$ (CHL ladder,
$\cC_N = D^b\Coh^{g_N}(\mathrm{K3} \times E)$), or
\item Scope (B): $N \in \{1, 2, 3, 4, 5, 6, 1/2, 3/2\}$
(Gritsenko--Cl\'ery $8$-form index).
\end{itemize}
Let $\PhiFA_3(\cC_N) \in E_3$-$\HolFA(X_N)$ be the Stage-$1$
$E_3$-factorisation algebra (Agent $1$ canonicity), $B^{E_3}$ the
Ayala--Francis Koszul-bar functor, and $\STr_{\modop}$ the
Getzler--Kapranov modular-operadic supertrace of
Definition~\ref{def:mod-op-supertrace-master}. Then
\begin{equation}
\label{eq:master-theorem-final}
 \kBKM(\Phi_N) \;=\; \frac{c_N(0)}{2}
 \;=\; \STr_{\modop}\bigl[B^{E_3}\bigl(\PhiFA_3(\cC_N)\bigr)\bigr].
\end{equation}
The identity is:
\begin{enumerate}[label=\textup{(\roman*)}]
\item canonical up to contractible choice (Heal~\ref{heal:three-path-qiso},
\ref{heal:scheme-independence});
\item $\GRT_1(\Q)$-invariant (Heal~\ref{heal:grt1-invariance});
\item Stage-$2$-invariant (Heal~\ref{heal:stage2-invariance});
\item compatible across the Scope (A) ladder and Scope (B) index
(Heal~\ref{heal:n-dependence});
\item consistent with the metaplectic and double-metaplectic covers
at Scope (B) half- and quarter-integer indices
(Heal~\ref{heal:metaplectic-scaling});
\item three-path--checkable via Hodge, BV, and MO paths
(Heal~\ref{heal:three-path-qiso}).
\end{enumerate}
\end{theorem}

\begin{proof}[Proof]
By Heal~\ref{heal:pairing-normalised}, the
modular-operadic pairing $\bigl\langle B^{E_3}(\PhiFA_3(\cC_N)),
[\Mbar_{2, 0}]\bigr\rangle^{\mathrm{mod, Hodge}}_{2, 0} = 1$
universally across $\cC_N$. By
Heal~\ref{heal:n-dependence}, the $N$-dependence enters via the
Borcherds-lift input: $\mathrm{Coeff}_{e^{2\pi i 0}}[\Phi_N] = c_N(0)$.
Combining through
Definition~\ref{def:mod-op-supertrace-master},
\[
 \STr_{\modop}\bigl[B^{E_3}(\PhiFA_3(\cC_N))\bigr] \;=\;
 1 \cdot \frac{c_N(0)}{2} \;=\; \frac{c_N(0)}{2} \;=\;
 \kBKM(\Phi_N)
\]
by Borcherds $1998$ Theorem~$13.3$. (i) follows from
Heal~\ref{heal:three-path-qiso} and
Heal~\ref{heal:scheme-independence};
(ii) from Heal~\ref{heal:grt1-invariance};
(iii) from Heal~\ref{heal:stage2-invariance};
(iv) from Heal~\ref{heal:n-dependence};
(v) from Heal~\ref{heal:metaplectic-scaling};
(vi) from Heal~\ref{heal:three-path-qiso}.
\end{proof}

% =====================================================================
\section{Worked examples}
\label{sec:worked-examples}
% =====================================================================

\subsection{$\cC_N = \Perf(\C^3)$ as base case}

\begin{proposition}[Base case: $\C^3$ produces $\kBKM = 0$]
\label{prop:base-case-c3}
For $\cC_0 = \Perf(\C^3)$, the Stage-$1$ lift $\PhiFA_3(\cC_0) \simeq
\mathrm{CoHA}(\C^3) \simeq Y^+(\widehat{\mathfrak{gl}}_1)$
(Schiffmann--Vasserot $2013$); the modular-operadic supertrace
evaluates to $0$ on $\mathrm{Sh}_{\mathrm{O}(2, 3)}$ because
$\C^3$ carries no Heegner divisor (non-compact host, $c_0(0) = 0$ by
convention). This is the Scope (B) $N = 3/2$ degenerate entry where
$\kBKM = 0$.
\end{proposition}

\subsection{$\cC_1 = D^b\Coh(\mathrm{K3} \times E)$: the flagship case}

\begin{proposition}[$\kBKM(\fg_{\Delta_5}) = 5$ via the master
identity]
\label{prop:k3xe-master-identity}
For $\cC_1 = D^b\Coh(\mathrm{K3} \times E)$, the master identity
at $N = 1$ reads
\[
 \kBKM(\Delta_5) \;=\; c_1(0)/2 \;=\; 10/2 \;=\; 5
 \;=\; \STr_{\modop}\bigl[B^{E_3}(\PhiFA_3(D^b\Coh(\mathrm{K3} \times E)))\bigr].
\]
The Borcherds input is the weight-$0$ index-$1$ weak Jacobi form
$\phi_{0, 1}(z, \tau) = \frac{12}{\theta_1^2(z, \tau)}[\eta(\tau)^{-2}
\cdot P^{\mathrm{ell}}(z, \tau)]$ where $P^{\mathrm{ell}}$ is the
K3 elliptic genus; its Borcherds lift is Gritsenko's paramodular cusp
form $\Delta_5$ of weight $5$ on $\Sp_4(\Z)$. The zeroth Fourier
coefficient of the input principal part is $c_1(0) = 10$ (the
$q$-constant term of $\phi_{0, 1}$ at $z = 0$ is $10$; see
Gritsenko--Nikulin $1995$ Part~II, Theorem~$2.1$; Lorgat $2020$
arXiv:$2004.09030$ Proposition~$3.2$).
\end{proposition}

\begin{proof}
Apply the master Theorem~\ref{thm:bl-master-kappa-bkm-chain-level-wave}
at $N = 1$, Scope (A) or Scope (B) (they coincide). The
modular-operadic pairing evaluates to $1$ on the
$\PhiFA_3(D^b\Coh(\mathrm{K3} \times E))$ Stage-$1$ lift because
$\mathrm{K3} \times E$ is a compact smooth CY$_3$ with cyclic
$A_\infty$-structure of parity $+3$ (Agent $14$ Heal~A.$3$). The
Borcherds input's zeroth Fourier coefficient is $10$ by Lorgat
$2020$ Proposition~$3.2$. The master identity gives $\kBKM =
10/2 = 5$.
\end{proof}

\begin{remark}[Hodge path cross-check on $\mathrm{K3} \times E$]
The Hodge path numeric is $\kch^{\mathrm{Hodge}}(\cC_1) =
\sum_{p, q} (-1)^{p + q} h^{p, q}(\mathrm{K3} \times E;
\Phi^*_{L_1})$; via Kapranov $L_\infty$ on
$T_{\mathrm{K3} \times E}[-1]$ (Agent $6$ three Atiyah cocycles
decomposition), the three cocycles contribute:
\begin{itemize}
\item $\At(T_{\mathrm{K3} \times E})$ contributes $0$ because
$c_1(T_{\mathrm{K3} \times E}) = 0$ in $H^{1, 1}$ (CY condition);
\item tree-level Yukawa $Y_3 \in H^{0, 3}(\mathrm{K3} \times E) = 0$
because $h^{0, 3}(\mathrm{K3} \times E) = 1$ but the $Y_3$ value on
the product CY is $0$ by K\"unneth ($Y_3(\mathrm{K3}) = 0$ trivially,
$Y_3(E)$ is purely $1$-dimensional, their K\"unneth product
vanishes in the tree-level bidegree);
\item one-loop BCOV curving $\alpha_{\BCOV} = (\chi_{\mathrm{top}}(
\mathrm{K3} \times E)/24) \cdot [\Omega_{\mathrm{K3} \times E}]^{0,
1} = 0$ by $\chi_{\mathrm{top}}(\mathrm{K3} \times E) = 24 \cdot 0 =
0$.
\end{itemize}
The Hodge-path raw numeric is $0$; this matches the Hodge supertrace
$\kch(\cC_1) = \Xi(\mathrm{K3} \times E) = 0$. The master identity's
value $5$ comes from the modular-operadic supertrace against the
Borcherds lift $\Delta_5$, not from the raw Hodge supertrace on
$\mathrm{K3} \times E$ alone: it is the Hodge supertrace on the
\emph{Heegner-divisor-paired} $\mathrm{K3} \times E$, which is
the Borcherds lift's zeroth Fourier coefficient.
\end{remark}

\subsection{The CHL ladder $N \in \{2, 3, 4, 6\}$}

\begin{proposition}[Master identity across the CHL ladder]
\label{prop:master-chl-ladder}
At Scope (A) $N \in \{2, 3, 4, 6\}$:
\[
 \kBKM(\Phi_N) \;=\; c_N(0)/2
 \;=\; (8, 6, 4, 2) / 2
 \;=\; (4, 3, 2, 1)
\]
at $N = (2, 3, 4, 6)$, respectively. The CY$_3$ host is the CHL
orbifold $Z_N = (\mathrm{K3} \times E)/\langle g_N \rangle$ with
$g_N \in M_{24}$ of order $N$ (Chaudhuri--Hockney--Lykken $1995$;
Govindarajan--Krishna $2010$ Table~1). The Borcherds input is the
Gritsenko lift of the $g_N$-twined K3 elliptic genus of the
relevant paramodular level (Gritsenko--Nikulin $1995$ Part~II,
Theorem~$2.1$).
\end{proposition}

\subsection{Scope (B): Gritsenko--Cl\'ery $8$-form family}

\begin{proposition}[Master identity across the
Gritsenko--Cl\'ery $8$-form family]
\label{prop:master-gc-8form}
At Scope (B) $N \in \{1, 2, 3, 4, 5, 6, 1/2, 3/2\}$:
\[
 \kBKM(\Phi_N) \;=\; c_N(0)/2
 \;=\; (10, 4, 2, 2, 1, 2, 1/2, 0) / 2
 \;=\; (5, 2, 1, 1, 1/2, 1, 1/4, 0),
\]
matching the eight-form catalogue of
Theorem~\ref{thm:eight-form-cy-host-catalogue} (target chapter).
The cover groups are:
\begin{itemize}
\item $\Sp_4(\Z)$ for $N \in \{1, 2, 3, 4, 6\}$ (integer weights);
\item $\Mp_4$ for $N = 5$ (half-integer weight via Shimura--Kohnen);
\item $\Mp_4$ for $N = 1/2$ (half-integer weight via Kohnen
double-cover);
\item $\widetilde{\Mp}_4$ for $N = 3/2$ (quarter-integer weight via
Bruinier--Funke; spin-double-cover).
\end{itemize}
\end{proposition}

% =====================================================================
\section{Cross-check with the three-path compatibility triangle}
\label{sec:three-path-cross-check}
% =====================================================================

For $N = 1$ (flagship case):
\begin{enumerate}[label=\textup{(\roman*)}]
\item Hodge path: $\kch^{\Hodge}(\cC_1; \Phi^*_{L_1}) = 5$,
computed as the Hodge supertrace of $\mathrm{K3} \times E$ twisted by
the Heegner-divisor line bundle of $\Phi_1 = \Delta_5$ (Bruinier
$2002$ Chern class formula).
\item BV path: $\kch^{\BV}(\cC_1) = 5$, computed as the Costello--Li--
Yagi one-loop partition function of $\hCS$ on $\mathrm{K3} \times E$
evaluated at the Heegner locus with Jacobi input $\phi_{0, 1}$
(Costello--Li $2016$ Proposition~$5.2$).
\item MO path: $\kch^{\MO}(\cC_1) = 5$, computed as the
Maulik--Okounkov $R$-matrix residue at the $N = 1$ Heegner wall on
$\mathrm{Sh}_{\mathrm{O}(2, 3)}$.
\end{enumerate}

The three paths agree numerically, confirming the three-path
compatibility triangle
Heal~\ref{heal:three-path-qiso}.

For $N = 6$ (Scope A boundary case):
\begin{enumerate}[label=\textup{(\roman*)}]
\item Hodge path: $\kch^{\Hodge}(\cC_6; \Phi^*_{L_6}) = 1$, computed
via $g_6$-twined Hodge supertrace of the $N = 6$ CHL quotient.
\item BV path: $\kch^{\BV}(\cC_6) = 1$, via the $g_6$-twisted Jacobi
input $\phi^{(g_6)}_{0, 1}$ in the Costello--Li one-loop formula.
\item MO path: $\kch^{\MO}(\cC_6) = 1$, via the $N = 6$ Heegner wall
residue.
\end{enumerate}

All three paths compute $1$, matching $\kBKM(\Phi_6) = c_6(0)/2 = 2/2
= 1$.

% =====================================================================
\section{Comparison with Scope (B) half-integer weight at $N = 5$}
\label{sec:half-integer-n5}
% =====================================================================

The Scope (B) $N = 5$ entry is the unique integer-index value with
half-integer weight $\kBKM = 1/2$, $c_5(0) = 1$. The Borcherds lift
$\Phi_5$ is a Siegel--Jacobi form on the metaplectic cover $\Mp_4 \to
\Sp_4$; its automorphy factor is the half-integer Maass character
$\nu_{1/2}$ (Kohnen $1982$).

\begin{proposition}[Master identity at $N = 5$ via metaplectic
scaling]
\label{prop:master-n5}
The master identity at $N = 5$ reads
\[
 \kBKM(\Phi_5) \;=\; c_5(0)/2 \;=\; 1/2 \;=\;
 \STr_{\modop}\bigl[B^{E_3}(\PhiFA_3(\cC_5))\bigr],
\]
with $\STr_{\modop}$ computed on the metaplectic cover
$\widetilde{\Mbar}_2$ via $[\widetilde{\Mbar}_{2, 0}] = 1/2 \cdot
[\Mbar_{2, 0}]$ (Kohnen $1982$); the pairing against the Borcherds
lift $\Phi_5 \in H^0(\widetilde{\Mbar}_2; \widetilde{\omega}^{1/2})$
gives $\kBKM = 1/2$.

\emph{Cross-check via the Hodge path.} The Hodge supertrace of the
$N = 5$ CY$_3$ host twisted by the Heegner line bundle
$\Phi^*_{L_5}$ is $1/2$; the half-integer arises because $\Phi_5$ is
a section of $\omega^{1/2}$, requiring a half-integer Chern character
computation (Bruinier $2002$ refined Heegner Chern-class formula).
\end{proposition}

% =====================================================================
\section{Quarter-integer weight at $N = 3/2$}
\label{sec:quarter-integer-n-3half}
% =====================================================================

\begin{proposition}[Master identity at $N = 3/2$ via double-metaplectic
scaling]
\label{prop:master-n-3half}
At $N = 3/2$ (Scope B) the Borcherds product is the degenerate
weight-$1/4$ form with $c_{3/2}(0) = 1/2$, $\kBKM = 1/4$. The
chain-level modular-operadic pairing reads
\[
 \STr_{\modop}\bigl[B^{E_3}(\PhiFA_3(\cC_{3/2}))\bigr]
 \;=\; \mathrm{Coeff}_{e^{2\pi i 0}}[\Phi_{3/2}] \cdot 1 \;=\; 1/4,
\]
on the double-metaplectic cover $\widetilde{\widetilde{\Mbar}}_2$,
where the fundamental class scales by $1/4$ per
Heal~\ref{heal:metaplectic-scaling}. The master identity reads
$\kBKM(\Phi_{3/2}) = 1/4$, matching the Bruinier--Funke $2004$ lift.
\end{proposition}

% =====================================================================
\section{Weight-$0$ degenerate terminal at $N = 8$}
\label{sec:terminal-n8}
% =====================================================================

\begin{proposition}[Master identity at $N = 8$ degenerate]
\label{prop:master-n8-degenerate}
At $N = 8$ (Scope B terminal) the Borcherds product $\Phi_8$ has
weight $0$: it is not a genuine Borcherds cusp form but the constant
$1$, with $c_8(0) = 0$, $\kBKM = 0$. The chain-level
modular-operadic pairing reads
\[
 \STr_{\modop}\bigl[B^{E_3}(\PhiFA_3(\cC_8))\bigr]
 \;=\; 0
\]
(the zeroth Fourier coefficient of the constant function $1$ is $0$
by the Borcherds-weight convention). The master identity reads
$\kBKM = 0$, the degenerate terminal fibre. This is the Scope (B)
$N = 8$ abelian-lattice entry.
\end{proposition}

% =====================================================================
\section{Summary of the master identity}
\label{sec:master-summary}
% =====================================================================

\begin{theorem}[Summary]
\label{thm:master-summary}
The chain-level master identity
$\kBKM(\Phi_N) = c_N(0)/2 = \STr_{\modop}[B^{E_3}(\PhiFA_3(\cC_N))]$
holds
\begin{enumerate}[label=\textup{(\roman*)}]
\item at Scope (A) across the CHL ladder $N \in \{1, 2, 3, 4, 6\}$,
with values $\kBKM = (5, 4, 3, 2, 1)$;
\item at Scope (B) across the Gritsenko--Cl\'ery $8$-form index
$N \in \{1, 2, 3, 4, 5, 6, 1/2, 3/2\}$, with values
$\kBKM = (5, 2, 1, 1, 1/2, 1, 1/4, 0)$;
\item coinciding at $N = 1$ on the Igusa--Gritsenko weight-$5$
Siegel paramodular cusp form $\Delta_5$;
\item three-path--checkable via Hodge, BV, and MO paths
(Section~\ref{sec:three-path-cross-check});
\item $\GRT_1(\Q)$-invariant on Stage-$1$ lifts;
\item Stage-$2$-invariant via the Mumford boundary-collapse
identification;
\item compatible with metaplectic and double-metaplectic covers at
half- and quarter-integer weights.
\end{enumerate}
\end{theorem}

\begin{remark}[Answer to Ghost~\ref{ghost:universal-supertrace}]
The value $c_N(0)/2$ is obtained as a supertrace (not a pairwise
evaluation): it is the modular-operadic trace of
$B^{E_3}(\PhiFA_3(\cC_N))$ against the Getzler--Kapranov chain-level
fundamental class $[\Mbar_{2, 0}]$ localised to the Borcherds-lift
generating series. The supertrace structure is load-bearing: the
pairing evaluates to $1$ universally, with the
$N$-dependence carried by the Borcherds-lift input; it is the
composition of these two factors that produces the scalar.
\end{remark}

% =====================================================================
\section{Open directions and chain-level frontier}
\label{sec:open-directions}
% =====================================================================

\begin{enumerate}[label=\textup{(\arabic*)}]
\item \emph{Genus $g \geq 3$.} The Schottky obstruction at $g \geq 4$
breaks the Torelli chain; the modular-operadic bridge at higher genus
requires the Schottky-locus version of the pairing, which is
currently conjectural (Agent $15$ Cycle III Heal~2).
\item \emph{Non-CHL paramodular hosts.} The Scope (B) values
$N \in \{5, 7, 8\}$ do not arise from a CHL quotient of
$\mathrm{K3} \times E$; the CY$_3$-functor-level status of
$\PhiFA_3(\cC_N)$ for these $N$ is conjectural. The master identity
is established for the Borcherds lift side; the CY$_3$-attached
Stage-$1$ lift is not canonical until a CY$_3$ host is fixed.
\item \emph{Quaternionic higher cover structure.} Beyond
$\widetilde{\widetilde{\Mp}}_4$, the Scheithauer $1999$ Nil-free
Borcherds lifts live on quaternionic covers; the chain-level
fundamental class on such covers is outside the current scope.
\item \emph{Higher-dimensional CY.} The master identity generalises
to CY$_d$ via the Getzler--Kapranov $E_d$-modular operad of
$(\overline{\mathcal M}_{g, n}, \overline{\mathcal A}_g,
\mathrm{Sh}_{\mathrm{O}(d - 1, d)})$ at $d = 5$ (Fake Monster
sibling, per memory). The $N = 1$ analogue is the Fake Monster
Borcherds lift $\Phi_{12}$ on $\mathrm{II}_{25, 1}$; the chain-level
modular-operadic construction is the direct generalisation, with
genus $= (d - 1)/2 = 2$ replaced by $(d - 1)/2$.
\item \emph{Moonshine sibling.} At $d = 3$, the Scope (A) $N = 1$
case ($\Delta_5$) is the BKM sibling of Igusa $\Phi_{10}$
(Gritsenko--Nikulin); the $d = 5$ sibling is the Fake Monster; at
$d = 2$, the $\Delta_5$ sibling is the Mathieu moonshine weight-$0$
form. The master identity organises these siblings as a
$d$-parametrised family of chain-level modular-operadic supertraces.
\end{enumerate}

% =====================================================================
\section{Primary-literature cross-reference}
\label{sec:primary-lit}
% =====================================================================

\noindent\textbf{Modular operad.}
\begin{itemize}
\item Getzler--Kapranov $1998$ \emph{Modular operads}, Compositio
Math.~$110$ (1998) 65--126 --- definition of modular operads,
boundary-gluing axioms, chain-level fundamental class.
\item Kapranov $1999$ \emph{Rozansky--Witten invariants via
Atiyah classes}, Compositio Math.~$115$ (1999) 71--113 ---
$H_\bullet(\Mbar_{0, n + 1}) = \Ger^{\mathrm{cyc}}_3(n)$
identification.
\end{itemize}

\noindent\textbf{$E_n$-operad formality.}
\begin{itemize}
\item Kontsevich $1999$ \emph{Operads and motives in deformation
quantisation}, Lett. Math. Phys.~$48$ --- graph integral
formality for $E_2$.
\item Tamarkin $2003$ \emph{Another proof of Kontsevich's
formality theorem} --- $E_2$-formality via associators.
\item Fresse $2011$ \emph{Koszul duality of $E_n$-operads},
Selecta Math.~$17$ --- $E_n$-Koszul self-duality up to shift.
\item Ayala--Francis $2014$ \emph{Poincar\'e/Koszul duality},
arXiv:$1409.2478$ --- FM compactification pairing.
\item Willwacher $2015$ \emph{Kontsevich's graph complex and
$\GRT_1$}, Invent. Math.~$200$ --- $\GRT_1(\Q)$-torsor structure.
\item Fresse--Willwacher $2020$ \emph{The intrinsic formality of
$E_n$-operads}, Acta Math.~$224$ --- $E_n$-formality for
$n \geq 2$.
\end{itemize}

\noindent\textbf{Borcherds lifts.}
\begin{itemize}
\item Borcherds $1995$ \emph{Automorphic forms on $O_{s + 2, 2}(\R)$
and infinite products}, Invent. Math.~$120$ (1995) 161--213 ---
singular-theta correspondence.
\item Borcherds $1998$ \emph{Automorphic forms with singularities on
Grassmannians}, Invent. Math.~$132$ (1998) 491--562 --- Theorem $13.3$
weight formula.
\item Gritsenko--Nikulin $1995$ \emph{Automorphic forms and Lorentzian
Kac--Moody algebras II}, arXiv:alg-geom/$9504006$ --- CHL ladder
Borcherds lifts.
\item Gritsenko $1999$ \emph{Arithmetical lifting and its
applications}, in: Algebraic Geometry Santa Cruz 1995 --- additive
lift construction.
\item Gritsenko--Cl\'ery $2013$ arXiv:$1305.4884$ --- $8$-form
Siegel-theta catalogue.
\item Bruinier $2002$ \emph{Borcherds products on $O(2, l)$ and
Chern classes of Heegner divisors}, Lecture Notes in Math.~$1780$
--- Heegner Chern class formula.
\item Lorgat $2020$ arXiv:$2004.09030$ \emph{Automorphic corrections
for the GKM superalgebra $\fg_{\Delta_5}$} --- explicit Borcherds
lift of $\phi_{0, 1}$ to $\Delta_5$.
\end{itemize}

\noindent\textbf{Shimura--Kohnen tower.}
\begin{itemize}
\item Shimura $1973$ \emph{Modular forms of half-integral weight},
Ann. Math.~$97$ --- half-integer modular forms.
\item Kohnen $1982$ \emph{Newforms of half-integral weight},
Math. Ann.~$260$ --- Kohnen double cover.
\item Bruinier--Funke $2004$ \emph{On two geometric theta-lifts},
Compositio Math.~$140$ --- metaplectic extensions.
\end{itemize}

\noindent\textbf{Chain-level curvature collapse.}
\begin{itemize}
\item Mumford $1983$ \emph{Towards enumerative geometry of the moduli
space of curves}, in: Arithmetic and Geometry II --- boundary
relation.
\item Janda--Pandharipande--Pixton--Zvonkine $2018$, J. Amer. Math.
Soc.~$31$ --- Pixton ideal relations on $R^\bullet(\Mbar_{g, n})$.
\end{itemize}

\noindent\textbf{Physics cross-check (BV, MO).}
\begin{itemize}
\item Costello--Li $2016$ arXiv:$1606.00365$ \emph{Quantum BCOV
theory on Calabi--Yau manifolds} --- one-loop Quantum Master
Equation.
\item Maulik--Okounkov $2014$ arXiv:$1211.1287$ \emph{Quantum
groups and quantum cohomology} --- $R$-matrix wall crossing.
\item Costello--Gaiotto--Yagi $2019$ \emph{Gauge theory and
integrability, IV} --- 5D $\hCS$ Yangian VOA.
\item Kapustin--Witten $2006$ --- Geometric Satake--Witten duality
framework.
\end{itemize}

% =====================================================================
\section*{Closing}
% =====================================================================

The master identity
\[
 \kBKM(\Phi_N) \;=\; c_N(0)/2 \;=\;
 \STr_{\modop}\bigl[B^{E_3}(\PhiFA_3(\cC_N))\bigr]
\]
is a chain-level theorem at Scope (A) $\cup$ Scope (B), with explicit
Koszul-duality control over Stage-$1$, explicit Stage-$2$-invariance
via the Mumford boundary collapse, explicit $\GRT_1(\Q)$-invariance
via Willwacher's graph-complex structure, and three-path
compatibility (Hodge, BV, MO). The flagship example $\cC_1 =
D^b\Coh(\mathrm{K3} \times E)$ recovers $\kBKM(\Delta_5) = 5$ via
Lorgat $2020$'s explicit construction; the extension to the metaplectic
and double-metaplectic covers recovers the Scope (B) half- and
quarter-integer values. The identity organises the
$\mathrm{BKM}$-modular-characteristic universally as a chain-level
invariant of the CY$_3$-attached Stage-$1$ $E_3$-algebra paired with
the Borcherds-lift generating series. The proof is synthetic:
Getzler--Kapranov modular operad, Ayala--Francis Koszul duality,
Torelli--Kuga--Satake uniformisation, Bruinier Heegner Chern-class
arithmetic, and the Costello--Li--Yagi one-loop BV anomaly cancellation
interlock to force the identity.

\end{document}
